% !TeX root = ../thuthesis-example.tex

\begin{denotation}[3cm]
  \item[$n$] 节点数量
  \item[$p_i$] 第$i$个节点的节点数量，$p_i \geq 0$并且$\sum_{i=1}^{n}p_i=1$
  \item[$\textbf{x}_t^{i}$] 第$i$个节点在第$t$轮迭代时的局部模型参数
  \item[$\textbf{x}_t$] $\textbf{x}_t=\sum_{i=1}^{n}p_i\textbf{x}_t^{i}$
  \item[$f_i(\textbf{x})$] 第$i$个节点的优化目标函数
  \item[$f(\textbf{x})$] $f(\textbf{x}):=\sum_{i=1}^{n}p_if_i(\textbf{x})$
  \item[$\nabla f(\textbf{x})$] $f(\textbf{x})$的梯度
  \item[$L$] $f$ 是 $L$光滑的
  \item[$\mu$] $f$ 是 $\mu$凸性的
  \item[$\delta$] 压缩率
  \item[$\textbf{e}_t^{i}$] 局部误差项
  \item[$\textbf{e}_t$]  $\textbf{e}_t=\sum_{i=1}^{n}p_i\textbf{e}_t^{i}$
  \item[$\xi^i$] 第$i$个节点的梯度噪音
  \item[$\zeta_i$] 衡量从节点$i$的局部分布到全局分布的距离参数
  \item[$\Bar{\zeta}$] 衡量全局数据异构性的参数
  \item[$M,\sigma^2$] 两个约束梯度噪声的常数
  \item[$G_t$] $G_t:=\textbf{E}\lVert \nabla f(x_t)\rVert^2$，在第$t$次迭代中对目标函数$f$的数学期望
  \item[$E_t$] 第$t$次迭代中总误差的数学期望
  \iffalse
  \item[DML] 聚苯基对称三嗪模型化合物的质子化产物
  \item[HMPPQ] 聚苯基喹噁啉模型化合物的质子化产物
  \item[PDT] 热分解温度
  \item[HPLC] 高效液相色谱（High Performance Liquid Chromatography）
  \item[HPCE] 高效毛细管电泳色谱（High Performance Capillary lectrophoresis）
  \item[LC-MS] 液相色谱-质谱联用（Liquid chromatography-Mass Spectrum）
  \item[TIC] 总离子浓度（Total Ion Content）
  \item[\textit{ab initio}] 基于第一原理的量子化学计算方法，常称从头算法
  \item[DFT] 密度泛函理论（Density Functional Theory）
  \item[$E_a$] 化学反应的活化能（Activation Energy）
  \item[ZPE] 零点振动能（Zero Vibration Energy）
  \item[PES] 势能面（Potential Energy Surface）
  \item[TS] 过渡态（Transition State）
  \item[TST] 过渡态理论（Transition State Theory）
  \item[$\increment G^\neq$] 活化自由能（Activation Free Energy）
  \item[$\kappa$] 传输系数（Transmission Coefficient）
  \item[IRC] 内禀反应坐标（Intrinsic Reaction Coordinates）
  \item[$\nu_i$] 虚频（Imaginary Frequency）
  \item[ONIOM] 分层算法（Our own N-layered Integrated molecular Orbital and molecular Mechanics）
  \item[SCF] 自洽场（Self-Consistent Field）
  \item[SCRF] 自洽反应场（Self-Consistent Reaction Field）
  \fi
\end{denotation}



% 也可以使用 nomencl 宏包，需要在导言区
% \usepackage{nomencl}
% \makenomenclature

% 在这里输出符号说明
% \printnomenclature[3cm]

% 在正文中的任意为都可以标题
% \nomenclature{PI}{聚酰亚胺}
% \nomenclature{MPI}{聚酰亚胺模型化合物，N-苯基邻苯酰亚胺}
% \nomenclature{PBI}{聚苯并咪唑}
% \nomenclature{MPBI}{聚苯并咪唑模型化合物，N-苯基苯并咪唑}
% \nomenclature{PY}{聚吡咙}
% \nomenclature{PMDA-BDA}{均苯四酸二酐与联苯四胺合成的聚吡咙薄膜}
% \nomenclature{MPY}{聚吡咙模型化合物}
% \nomenclature{As-PPT}{聚苯基不对称三嗪}
% \nomenclature{MAsPPT}{聚苯基不对称三嗪单模型化合物，3,5,6-三苯基-1,2,4-三嗪}
% \nomenclature{DMAsPPT}{聚苯基不对称三嗪双模型化合物（水解实验模型化合物）}
% \nomenclature{S-PPT}{聚苯基对称三嗪}
% \nomenclature{MSPPT}{聚苯基对称三嗪模型化合物，2,4,6-三苯基-1,3,5-三嗪}
% \nomenclature{PPQ}{聚苯基喹噁啉}
% \nomenclature{MPPQ}{聚苯基喹噁啉模型化合物，3,4-二苯基苯并二嗪}
% \nomenclature{HMPI}{聚酰亚胺模型化合物的质子化产物}
% \nomenclature{HMPY}{聚吡咙模型化合物的质子化产物}
% \nomenclature{HMPBI}{聚苯并咪唑模型化合物的质子化产物}
% \nomenclature{HMAsPPT}{聚苯基不对称三嗪模型化合物的质子化产物}
% \nomenclature{HMSPPT}{聚苯基对称三嗪模型化合物的质子化产物}
% \nomenclature{HMPPQ}{聚苯基喹噁啉模型化合物的质子化产物}
% \nomenclature{PDT}{热分解温度}
% \nomenclature{HPLC}{高效液相色谱（High Performance Liquid Chromatography）}
% \nomenclature{HPCE}{高效毛细管电泳色谱（High Performance Capillary lectrophoresis）}
% \nomenclature{LC-MS}{液相色谱-质谱联用（Liquid chromatography-Mass Spectrum）}
% \nomenclature{TIC}{总离子浓度（Total Ion Content）}
% \nomenclature{\textit{ab initio}}{基于第一原理的量子化学计算方法，常称从头算法}
% \nomenclature{DFT}{密度泛函理论（Density Functional Theory）}
% \nomenclature{$E_a$}{化学反应的活化能（Activation Energy）}
% \nomenclature{ZPE}{零点振动能（Zero Vibration Energy）}
% \nomenclature{PES}{势能面（Potential Energy Surface）}
% \nomenclature{TS}{过渡态（Transition State）}
% \nomenclature{TST}{过渡态理论（Transition State Theory）}
% \nomenclature{$\increment G^\neq$}{活化自由能（Activation Free Energy）}
% \nomenclature{$\kappa$}{传输系数（Transmission Coefficient）}
% \nomenclature{IRC}{内禀反应坐标（Intrinsic Reaction Coordinates）}
% \nomenclature{$\nu_i$}{虚频（Imaginary Frequency）}
% \nomenclature{ONIOM}{分层算法（Our own N-layered Integrated molecular Orbital and molecular Mechanics）}
% \nomenclature{SCF}{自洽场（Self-Consistent Field）}
% \nomenclature{SCRF}{自洽反应场（Self-Consistent Reaction Field）}
